\documentclass[10pt,a4paper]{amsart}
\usepackage{amsmath,amssymb,amsthm,graphicx}
% Retained minimal preamble for the Stacks Project chapter pione.tex (Stacks chapter, pre-delegation retained fragment for stacks-pione).
% Verbatim excerpt of the chapter preamble (sources/stacks-spaces-descent/preamble.tex): the reference/slogan/history
% environment block (native lines 7-12), Xy-pic setup (native lines 14-19) and the native theorem declarations with
% the mathematical macros (native lines 154-217). Copied mechanically; nothing added, removed or redefined. The
% wrapper class and the amsmath/amssymb/amsthm/graphicx base packages are supplied by the assembler, so the
% class/xr-hyper/hyperref/externaldocument part of the native preamble is omitted.
% For dealing with references we use the comment environment
\usepackage{verbatim}
\newenvironment{reference}{\comment}{\endcomment}
%\newenvironment{reference}{}{}
\newenvironment{slogan}{\comment}{\endcomment}
\newenvironment{history}{\comment}{\endcomment}
% For commutative diagrams we use Xy-pic
\usepackage[all]{xy}

% We use 2cell for 2-commutative diagrams.
\xyoption{2cell}
\UseAllTwocells
% Theorem environments.
%
\theoremstyle{plain}
\newtheorem{theorem}[subsection]{Theorem}
\newtheorem{proposition}[subsection]{Proposition}
\newtheorem{lemma}[subsection]{Lemma}

\theoremstyle{definition}
\newtheorem{definition}[subsection]{Definition}
\newtheorem{example}[subsection]{Example}
\newtheorem{exercise}[subsection]{Exercise}
\newtheorem{situation}[subsection]{Situation}

\theoremstyle{remark}
\newtheorem{remark}[subsection]{Remark}
\newtheorem{remarks}[subsection]{Remarks}

\numberwithin{equation}{subsection}

% Macros
%
\def\lim{\mathop{\mathrm{lim}}\nolimits}
\def\colim{\mathop{\mathrm{colim}}\nolimits}
\def\Spec{\mathop{\mathrm{Spec}}}
\def\Hom{\mathop{\mathrm{Hom}}\nolimits}
\def\Ext{\mathop{\mathrm{Ext}}\nolimits}
\def\SheafHom{\mathop{\mathcal{H}\!\mathit{om}}\nolimits}
\def\SheafExt{\mathop{\mathcal{E}\!\mathit{xt}}\nolimits}
\def\Sch{\mathit{Sch}}
\def\Mor{\mathop{\mathrm{Mor}}\nolimits}
\def\Ob{\mathop{\mathrm{Ob}}\nolimits}
\def\Sh{\mathop{\mathit{Sh}}\nolimits}
\def\NL{\mathop{N\!L}\nolimits}
\def\CH{\mathop{\mathrm{CH}}\nolimits}
\def\proetale{{pro\text{-}\acute{e}tale}}
\def\etale{{\acute{e}tale}}
\def\QCoh{\mathit{QCoh}}
\def\Ker{\mathop{\mathrm{Ker}}}
\def\Im{\mathop{\mathrm{Im}}}
\def\Coker{\mathop{\mathrm{Coker}}}
\def\Coim{\mathop{\mathrm{Coim}}}

% Boxtimes
%
\DeclareMathSymbol{\boxtimes}{\mathbin}{AMSa}{"02}

%
% Macros for moduli stacks/spaces
%
\def\QCohstack{\mathcal{QC}\!\mathit{oh}}
\def\Cohstack{\mathcal{C}\!\mathit{oh}}
\def\Spacesstack{\mathcal{S}\!\mathit{paces}}
\def\Quotfunctor{\mathrm{Quot}}
\def\Hilbfunctor{\mathrm{Hilb}}
\def\Curvesstack{\mathcal{C}\!\mathit{urves}}
\def\Polarizedstack{\mathcal{P}\!\mathit{olarized}}
\def\Complexesstack{\mathcal{C}\!\mathit{omplexes}}
% \Pic is the operator that assigns to X its picard group, usage \Pic(X)
% \Picardstack_{X/B} denotes the Picard stack of X over B
% \Picardfunctor_{X/B} denotes the Picard functor of X over B
\def\Pic{\mathop{\mathrm{Pic}}\nolimits}
\def\Picardstack{\mathcal{P}\!\mathit{ic}}
\def\Picardfunctor{\mathrm{Pic}}
\def\Deformationcategory{\mathcal{D}\!\mathit{ef}}
\usepackage{cleveref}
\crefname{originallemma}{original lemma}{original lemmas}
\Crefname{originallemma}{Original lemma}{Original lemmas}
\makeatletter
\newlabel{algebra-lemma-Sk-goes-up}{{0339}{source}{}{sourcepointer-1}{}}
\newlabel{algebra-lemma-Sk-goes-up@cref}{{[originallemma][0][]0339}{[1][0][]source}{}{}{}}
\makeatother
\crefname{originallemma}{original lemma}{original lemmas}
\Crefname{originallemma}{Original lemma}{Original lemmas}
\makeatletter
\newlabel{algebraization-lemma-fully-faithful-simple-two}{{0EKT}{source}{}{sourcepointer-2}{}}
\newlabel{algebraization-lemma-fully-faithful-simple-two@cref}{{[originallemma][0][]0EKT}{[1][0][]source}{}{}{}}
\makeatother
\crefname{originallemma}{original lemma}{original lemmas}
\Crefname{originallemma}{Original lemma}{Original lemmas}
\makeatletter
\newlabel{dualizing-lemma-universally-catenary}{{0A80}{source}{}{sourcepointer-3}{}}
\newlabel{dualizing-lemma-universally-catenary@cref}{{[originallemma][0][]0A80}{[1][0][]source}{}{}{}}
\makeatother
\crefname{originallemma}{original lemma}{original lemmas}
\Crefname{originallemma}{Original lemma}{Original lemmas}
\makeatletter
\newlabel{dualizing-section-formal-fibres}{{0BJM}{source}{}{sourcepointer-4}{}}
\newlabel{dualizing-section-formal-fibres@cref}{{[originallemma][0][]0BJM}{[1][0][]source}{}{}{}}
\makeatother
\crefname{originallemma}{original lemma}{original lemmas}
\Crefname{originallemma}{Original lemma}{Original lemmas}
\makeatletter
\newlabel{lemma-fully-faithful-henselian-completion}{{0BLI}{source}{}{sourcepointer-5}{}}
\newlabel{lemma-fully-faithful-henselian-completion@cref}{{[originallemma][0][]0BLI}{[1][0][]source}{}{}{}}
\makeatother
\crefname{originallemma}{original lemma}{original lemmas}
\Crefname{originallemma}{Original lemma}{Original lemmas}
\makeatletter
\newlabel{lemma-restriction-fully-faithful}{{0EL8}{source}{}{sourcepointer-6}{}}
\newlabel{lemma-restriction-fully-faithful@cref}{{[originallemma][0][]0EL8}{[1][0][]source}{}{}{}}
\makeatother
\crefname{originallemma}{original lemma}{original lemmas}
\Crefname{originallemma}{Original lemma}{Original lemmas}
\makeatletter
\newlabel{local-cohomology-lemma-change-completion}{{0EHW}{source}{}{sourcepointer-7}{}}
\newlabel{local-cohomology-lemma-change-completion@cref}{{[originallemma][0][]0EHW}{[1][0][]source}{}{}{}}
\makeatother
\crefname{originallemma}{original lemma}{original lemmas}
\Crefname{originallemma}{Original lemma}{Original lemmas}
\makeatletter
\newlabel{more-algebra-proposition-ratliff}{{0AW6}{source}{}{sourcepointer-8}{}}
\newlabel{more-algebra-proposition-ratliff@cref}{{[originallemma][0][]0AW6}{[1][0][]source}{}{}{}}
\makeatother
\crefname{originallemma}{original lemma}{original lemmas}
\Crefname{originallemma}{Original lemma}{Original lemmas}
\makeatletter
\newlabel{more-algebra-section-properties-formal-fibres}{{0BIR}{source}{}{sourcepointer-9}{}}
\newlabel{more-algebra-section-properties-formal-fibres@cref}{{[originallemma][0][]0BIR}{[1][0][]source}{}{}{}}
\makeatother
\crefname{originallemma}{original lemma}{original lemmas}
\Crefname{originallemma}{Original lemma}{Original lemmas}
\makeatletter
\newlabel{situation-local-lefschetz}{{0BLF}{source}{}{sourcepointer-10}{}}
\newlabel{situation-local-lefschetz@cref}{{[originallemma][0][]0BLF}{[1][0][]source}{}{}{}}
\makeatother
\title{Local Lefschetz full faithfulness for the etale punctured neighbourhood under a dualizing complex, henselian pair, and either an S-2 plus three-dimensional-component condition or a depth-versus-dimension inequality strictly above two.}
\author{The Stacks Project Authors}
\date{stacks-project:pione.tex:local-snapshot-20261012; GNU FDL 1.2 or later}
\begin{document}
\maketitle
\noindent\textit{Editorial note.} Statement span is lines 4816-4834 of sources/stacks-pione/pinned-original-source.tex; proof span is lines 4836-4860, the complete author proof (23 body lines).\par\medskip
\noindent\textit{Editorial note.} Navigation stubs: 10 cited-but-not-included labels are supplied as non-load-bearing pointers with their pinned tags where the snapshot tag index has them.\par\medskip
\noindent\textit{Editorial note.} extra\_packages is empty on purpose: the retained preamble already loads xy with options and a second declaration clashes in the assembler wrapper.\par\medskip
\noindent\textit{Editorial note.} Text is retained verbatim, including the author's own typos and placeholders.\par\medskip
\noindent\textit{Editorial note.} Pool stacks-pione wave 34d.\par\medskip
\noindent\textit{Editorial note.} rqual sv-mismatch fix 2026-10-14: source\_version aligned to SOURCE.json (semicolon+git form, intake authority; package evidence.json already correct); revision bumped for re-assembly per publication rotation (takeover-line defer ticket, relayed by the s15 line).\par\medskip
\noindent\textit{Allowed standard prerequisites (editorial).}\begin{itemize}
\item Explicit prerequisite boundary. See why-field; every pointer named here occurs verbatim in the retained statement/proof text.
\end{itemize}\medskip
\ifdefined\hypertarget\hypertarget{sourcepointer-1}{}\fi
\noindent\textit{Source reference.} Original lemma 0339: sources/stacks-pione/pinned-original-source.tex, source lines 4860--4860. Cited by the retained statement or proof but not part of the collected spans; supplied only as a navigation pointer. The argument does not depend on its content.\par\medskip
\ifdefined\hypertarget\hypertarget{sourcepointer-2}{}\fi
\noindent\textit{Source reference.} Original lemma 0EKT: sources/stacks-pione/pinned-original-source.tex, source lines 4860--4860. Cited by the retained statement or proof but not part of the collected spans; supplied only as a navigation pointer. The argument does not depend on its content.\par\medskip
\ifdefined\hypertarget\hypertarget{sourcepointer-3}{}\fi
\noindent\textit{Source reference.} Original lemma 0A80: sources/stacks-pione/pinned-original-source.tex, source lines 4860--4860. Cited by the retained statement or proof but not part of the collected spans; supplied only as a navigation pointer. The argument does not depend on its content.\par\medskip
\ifdefined\hypertarget\hypertarget{sourcepointer-4}{}\fi
\noindent\textit{Source reference.} Original lemma 0BJM: sources/stacks-pione/pinned-original-source.tex, source lines 4860--4860. Cited by the retained statement or proof but not part of the collected spans; supplied only as a navigation pointer. The argument does not depend on its content.\par\medskip
\ifdefined\hypertarget\hypertarget{sourcepointer-5}{}\fi
\noindent\textit{Source reference.} Original lemma 0BLI: sources/stacks-pione/pinned-original-source.tex, source lines 4733--4741. Cited by the retained statement or proof but not part of the collected spans; supplied only as a navigation pointer. The argument does not depend on its content.\par\medskip
\ifdefined\hypertarget\hypertarget{sourcepointer-6}{}\fi
\noindent\textit{Source reference.} Original lemma 0EL8: sources/stacks-pione/pinned-original-source.tex, source lines 4175--4189. Cited by the retained statement or proof but not part of the collected spans; supplied only as a navigation pointer. The argument does not depend on its content.\par\medskip
\ifdefined\hypertarget\hypertarget{sourcepointer-7}{}\fi
\noindent\textit{Source reference.} Original lemma 0EHW: sources/stacks-pione/pinned-original-source.tex, source lines 4860--4860. Cited by the retained statement or proof but not part of the collected spans; supplied only as a navigation pointer. The argument does not depend on its content.\par\medskip
\ifdefined\hypertarget\hypertarget{sourcepointer-8}{}\fi
\noindent\textit{Source reference.} Original lemma 0AW6: sources/stacks-pione/pinned-original-source.tex, source lines 4860--4860. Cited by the retained statement or proof but not part of the collected spans; supplied only as a navigation pointer. The argument does not depend on its content.\par\medskip
\ifdefined\hypertarget\hypertarget{sourcepointer-9}{}\fi
\noindent\textit{Source reference.} Original lemma 0BIR: sources/stacks-pione/pinned-original-source.tex, source lines 4860--4860. Cited by the retained statement or proof but not part of the collected spans; supplied only as a navigation pointer. The argument does not depend on its content.\par\medskip
\ifdefined\hypertarget\hypertarget{sourcepointer-10}{}\fi
\noindent\textit{Source reference.} Original lemma 0BLF: sources/stacks-pione/pinned-original-source.tex, source lines 4651--4657. Cited by the retained statement or proof but not part of the collected spans; supplied only as a navigation pointer. The argument does not depend on its content.\par\medskip
\begin{lemma}
\label{lemma-fully-faithful-simple}
In Situation \ref{situation-local-lefschetz}. Assume
\begin{enumerate}
\item[(a)] $A$ has a dualizing complex,
\item[(b)] the pair $(A, (f))$ is henselian,
\item[(c)] one of the following is true
\begin{enumerate}
\item[(i)] $A_f$ is $(S_2)$ and every irreducible component of $X$
not contained in $X_0$ has dimension $\geq 3$, or
\item[(ii)] for every prime
$\mathfrak p \subset A$, $f \not \in \mathfrak p$ we have
$\text{depth}(A_\mathfrak p) + \dim(A/\mathfrak p) > 2$.
\end{enumerate}
\end{enumerate}
Then the restriction functor
$\textit{F\'Et}_U \longrightarrow \textit{F\'Et}_{U_0}$
is fully faithful.
\end{lemma}

\begin{proof}
Let $A'$ be the $\mathfrak m$-adic completion of $A$. We will show that
the hypotheses remain true for $A'$. This is clear for conditions
(a) and (b). Condition (c)(ii) is preserved by
Local Cohomology, Lemma \ref{local-cohomology-lemma-change-completion}.
Next, assume (c)(i) holds. Since $A$ is universally catenary
(Dualizing Complexes, Lemma \ref{dualizing-lemma-universally-catenary})
we see that every irreducible component of $\Spec(A')$ not contained in $V(f)$
has dimension $\geq 3$, see
More on Algebra, Proposition \ref{more-algebra-proposition-ratliff}.
Since $A \to A'$ is flat with Gorenstein fibres,
the condition that $A_f$ is $(S_2)$ implies that $A'_f$ is $(S_2)$.
References used:
Dualizing Complexes, Section \ref{dualizing-section-formal-fibres},
More on Algebra, Section \ref{more-algebra-section-properties-formal-fibres},
and Algebra, Lemma \ref{algebra-lemma-Sk-goes-up}.
Thus by Lemma \ref{lemma-fully-faithful-henselian-completion}
we may assume that $A$ is a Noetherian complete local ring.

\medskip\noindent
Assume $A$ is a complete local ring in addition to the other assumptions.
By Lemma \ref{lemma-restriction-fully-faithful} the result follows from
Algebraic and Formal Geometry, Lemma
\ref{algebraization-lemma-fully-faithful-simple-two}.
\end{proof}

\end{document}
